% !TEX root = ../../main.tex
\subsection{订单 Token 认证机制}
\label{subsec:token}

\subsubsection*{机制概述}

订单 Token 是用户进入订单通信通道前必须提交的认证凭证。Token 由服务器在订单创建时生成，并与订单绑定保存。用户发起 WebSocket 入房请求时，服务器对 Token 的合法性、归属关系和时效性进行三重验证，缺一不可。

该机制的核心设计思路是：Token 不是一个简单的随机字符串，而是将订单号、用户 PID、时间戳和随机数绑定在一起的 HMAC-SM3 认证码。攻击者在不知道订单认证密钥$K_{\mathrm{order}}$的情况下，无法伪造任何合法 Token，也无法从一个合法 Token 推导出其他订单的凭证。

\subsubsection*{Token 生成公式}

Token 的理想生成公式为：
\[
  \mathrm{Token} = \operatorname{HMAC\text{-}SM3}\!\left(K_{\mathrm{order}},\; \mathit{orderID} \| \mathrm{PID} \| \mathit{timestamp} \| \mathit{nonce}\right)
\]
其中：
\begin{itemize}[leftmargin=2em]
  \item$K_{\mathrm{order}}$为系统订单认证密钥，由服务端统一管理；
  \item$\mathit{orderID}$为随机 UUID 订单号，防止顺序枚举；
  \item$\mathrm{PID}$为发起订单的用户匿名身份标识；
  \item$\mathit{timestamp}$为 Token 生成时的 Unix 时间戳（秒级），用于有效期校验；
  \item$\mathit{nonce}$为随机数，进一步降低重放攻击风险；
  \item Token 生成后，$\langle\mathit{token}, \mathit{timestamp}\rangle$一并写入 \texttt{orders} 表供后续验证使用。
\end{itemize}

引入$\mathit{nonce}$的目的在于：即使攻击者能够在 Token 有效期内截获同一组（orderID、PID、timestamp）的请求，相同的时间戳配合不同的 nonce 仍会产生不同的 Token 值，从而进一步降低重放攻击成功率。

\subsubsection*{Token 验证流程}

用户提交 \texttt{join\_order} 请求时，服务器执行以下验证步骤：

% !TEX root = ../../../main.tex
% Token 三重验证流程图
\begin{figure}[H]
\centering
\begin{tikzpicture}[
  node distance = 0.55cm and 2.6cm,
  proc/.style     = {rectangle, draw, rounded corners=3pt, minimum width=3.4cm, minimum height=0.65cm,
                     font=\small, fill=blue!8, draw=blue!50!black},
  decision/.style = {diamond, draw, aspect=2.4, minimum width=3.4cm, minimum height=0.8cm,
                     font=\small, align=center, fill=orange!15, draw=orange!60!black, inner sep=2pt},
  term/.style     = {rectangle, draw, rounded corners=8pt, minimum width=3.2cm, minimum height=0.65cm,
                     font=\small, fill=green!14, draw=green!50!black},
  reject/.style   = {rectangle, draw, rounded corners=8pt, minimum width=2.4cm, minimum height=0.62cm,
                     font=\small, fill=red!12, draw=red!50!black},
  arr/.style      = {-Stealth, thick},
  lbl/.style      = {font=\small},
]

\node[proc]     (recv)    {接收 orderID + token};
\node[decision, below=0.75cm of recv]   (d_time)  {时间戳\\在有效期内？};
\node[reject,   right=2.6cm of d_time]  (r_time)  {拒绝：已过期};
\node[proc,     below=0.9cm of d_time]  (recalc)  {重算$\mathit{token}'{=}\mathrm{HMAC}$-$\mathrm{SM3}$};
\node[decision, below=0.8cm of recalc]  (d_eq)    {token'\\== token？};
\node[reject,   right=2.6cm of d_eq]   (r_eq)    {拒绝：签名不匹配};
\node[decision, below=0.9cm of d_eq]   (d_stat)  {订单状态\\为活跃？};
\node[reject,   right=2.6cm of d_stat] (r_stat)  {拒绝：订单已关闭};
\node[term,     below=0.9cm of d_stat] (pass)    {认证通过，加入订单房间};

\draw[arr] (recv)   -- (d_time);
\draw[arr] (d_time) -- node[lbl,left]{是} (recalc);
\draw[arr] (d_time) -- node[lbl,above]{否} (r_time);
\draw[arr] (recalc) -- (d_eq);
\draw[arr] (d_eq)   -- node[lbl,left]{是} (d_stat);
\draw[arr] (d_eq)   -- node[lbl,above]{否} (r_eq);
\draw[arr] (d_stat) -- node[lbl,left]{是} (pass);
\draw[arr] (d_stat) -- node[lbl,above]{否} (r_stat);

\end{tikzpicture}
\caption{订单 Token 三重验证流程图}
\label{fig:token_verify}
\end{figure}


\begin{enumerate}[label=\arabic*., leftmargin=2em]
  \item \textbf{有效期校验}：读取请求中的$\mathit{timestamp}$，与当前服务器时间比较，若超出有效窗口（默认 3600 秒），拒绝本次请求；
  \item \textbf{重新计算 Token}：以数据库中保存的$K_{\mathrm{order}}$、$\mathit{orderID}$、$\mathrm{PID}$、$\mathit{timestamp}$（及$\mathit{nonce}$）为输入，重新调用 HMAC-SM3 计算期望 Token；
  \item \textbf{比对验证}：将重新计算得到的期望 Token 与请求中提交的 Token 进行恒定时间比较，一致则通过，否则拒绝；
  \item \textbf{订单状态检查}：验证通过后进一步检查订单状态，确保订单处于合法状态（如未被取消），方可允许用户加入通信房间。
\end{enumerate}

\subsubsection*{安全性分析}

Token 机制提供了以下安全保障：

\begin{itemize}[leftmargin=2em]
  \item \textbf{抗伪造性}：Token 由 HMAC-SM3 生成，在$K_{\mathrm{order}}$保密的前提下，攻击者无法在不知道密钥的情况下构造合法 Token；
  \item \textbf{绑定性}：Token 将$\mathit{orderID}$、$\mathrm{PID}$、$\mathit{timestamp}$、$\mathit{nonce}$绑定为一个整体，任意字段被篡改均会导致验证失败；
  \item \textbf{时效性}：时间戳有效期机制限制 Token 的可用窗口，降低截获后长期利用的风险；
  \item \textbf{抗重放性}：$\mathit{nonce}$的引入保证即使相同的订单号和时间戳也会产生不同的 Token，进一步降低重放成功概率。
\end{itemize}
